% -----------------------------------------------------------------------------
% Section 2: Related Work
% Paper: "Silent Decisions Are Type Errors: Enforcing AI Accountability
%         via Linear Resource Types"
% Venue: IEEE Computer — Special Issue on AI Governance
% Deadline: 13 April 2026
% Authors: Daniel Lau Pereira Soares
% Rev 3: 25 March 2026 — ACID distinction added (§2.2)
% -----------------------------------------------------------------------------

\section{Related Work}
\label{sec:related_work}

The challenge of algorithmic accountability intersects three distinct
domains of computer science: infrastructure governance, runtime
verification, and type theory.
We situate our contribution in contrast to prevailing methodologies
in each domain.

% ----------------------------------------------------------------
\subsection{Policy-as-Code for Infrastructure Governance}
\label{sec:pac}

The industry standard for continuous compliance relies on
Policy-as-Code (PaC) frameworks such as Open Policy
Agent~\cite{opa2024}, Kyverno, and Checkov.
These tools excel at governing infrastructure artifacts---verifying
that a Kubernetes pod lacks privilege escalation, that a cloud
storage bucket is not publicly accessible, or that a Terraform plan
does not provision over-privileged IAM roles.
However, they are designed to evaluate \emph{declarative
configurations}, not \emph{dynamic runtime decisions} affecting
individual rights.

Consequently, evidence of compliance in these systems manifests as
asynchronous audit logs: records written after the decision logic has
executed, stored in a separate artifact, and subject to omission or
alteration if the logging pipeline fails.
Infrastructure PaC frameworks do not address the problem we target,
because they assume a world where the object of governance is a
configuration artifact rather than a transient decision act with
immediate legal consequences for a natural person.
In PaC, accountability is a post-hoc observable state---a property
that may or may not be present at audit time, leaving the system
structurally vulnerable to the silent-decision failure mode we
identify in Section~\ref{sec:intro}.

% ----------------------------------------------------------------
\subsection{Runtime Assertions and Audit Logging}
\label{sec:runtime}

To govern application-level decision logic, practitioners
typically rely on a combination of runtime assertions
(e.g., \texttt{assert(evidence != null)}) and structured
audit logging.
While functional under nominal conditions, these mechanisms are
procedural promises, not structural invariants.
Runtime assertions can be disabled in production builds
(e.g., by stripping \texttt{debug\_assertions} in Rust, or by
omitting \texttt{-ea} in the Java VM) to reduce overhead.
Audit logs can be truncated under load, dropped by a misconfigured
sink, or bypassed by a code path that was never reached during
testing.

Moreover, both mechanisms operate \emph{after} the decision logic
has executed.
They do not prevent the materialization of a verdict that lacks
non-repudiable evidence; they attempt to record its absence, if they
run at all.
This architectural decoupling separates the act of deciding from the
act of proving---a separation that is incompatible with the
structural accountability obligations imposed by LGPD Art.~18\S2 and
EU AI Act Art.~86, which require that evidence \emph{accompany} a
decision, not merely follow it.

A stronger form of runtime guarantee is offered by transactional
databases: an ACID transaction wrapping both the decision record and
the audit log ensures that either both are committed or neither is.
However, transactional atomicity is a tool the developer must
consciously choose to apply---wrapping decision and log within the
same \texttt{BEGIN}/\texttt{COMMIT} boundary.
Omitting that boundary is a logic error, not a type error: the
compiler accepts the program, the database commits the decision
without the log, and the silent-decision failure mode materialises
with no static diagnostic.
BTV closes exactly this gap: the compiler rejects any program that
constructs a \texttt{Verdict} without a causally-prior
\texttt{EvidenceToken}, regardless of whether the developer intended
to include the evidence or not.

% ----------------------------------------------------------------
\subsection{Linear Types and Resource Correctness}
\label{sec:linear_related}

The theoretical foundation for structural resource guarantees is
Linear Logic, introduced by Girard~\cite{girard1987linear} in 1987.
Girard's key insight---that propositions should be treated as
\emph{resources} that are consumed exactly once---was subsequently
adapted into programming language design by
Wadler~\cite{wadler1991linear}, who showed that linear types
eliminate a class of resource-management errors at compile time.
This paradigm has proven highly effective in enforcing memory safety
and preventing data races; it underpins the ownership and borrowing
model of the Rust programming language~\cite{jung2017rustbelt}.

However, prior applications of linear types focus exclusively on
\emph{resource} correctness---ensuring that memory is freed exactly
once, that file handles are not leaked, that channels are not
reused after a session ends.
To the best of our knowledge, no prior work has applied linear
resource types to enforce \emph{legal} correctness: the requirement
that a decision act cannot exist without a causally-bound,
non-repudiable evidence chain.
Bound-To-Verdict bridges this gap, extending the linear-type
discipline from resource management to regulatory accountability,
and demonstrating that the same compile-time guarantees that prevent
memory unsafety can be applied to prevent accountability unsafety.
